% =========================================================================
% บทที่ 2: พลศาสตร์ขั้นสูงและสมการการเคลื่อนที่แบบไม่เชิงเส้น
% =========================================================================

\chapter{พลศาสตร์ขั้นสูงและสมการการเคลื่อนที่แบบไม่เชิงเส้น}
\label{ch:chapter02}

\noindent{\large\bfseries\color{physicsblue} แบบจำลองแรงต้านอากาศไม่เชิงเส้น ความเร็วปลาย และระเบียบวิธีรูงเง-คุตตา}

\vspace{0.4cm}

\begin{equationsbox}{สมการการเคลื่อนที่ภายใต้แรงต้านอากาศ}
\begin{align}
m\frac{d\mathbf{v}}{dt} &= m\mathbf{g} - b\mathbf{v} - c v \mathbf{v}, \quad v_t = \sqrt{\frac{mg}{c}} = \sqrt{\frac{2mg}{\rho A C_D}} \\
v(t) &= v_t \tanh\left(\frac{gt}{v_t}\right), \quad y(t) = \frac{v_t^2}{g}\ln\left[\cosh\left(\frac{gt}{v_t}\right)\right]
\end{align}
\end{equationsbox}

\section{ทฤษฎีและกรอบแนวคิดสำคัญ}

ในฟิสิกส์เบื้องต้น การเคลื่อนที่ของวัตถุภายใต้สนามโน้มถ่วงมักละเลยแรงต้านของอากาศ แต่ในความเป็นจริงทางวิศวกรรมและปรากฏการณ์บรรยากาศ แรงต้านอากาศ (Aerodynamic Drag Force) เป็นแรงไม่เชิงเส้นที่ขึ้นอยู่กับอัตราเร็วของวัตถุ สำหรับอัตราเร็วต่ำ (Low Reynolds Number, $Re < 1$) แรงต้านจะแปรผันตรงกับความเร็วเชิงเส้น $\mathbf{F}_D = -b\mathbf{v}$ ตามกฎของสโตกส์ แต่สำหรับวัตถุขนาดใหญ่ที่มีอัตราเร็วสูงในบรรยากาศ ($10^3 < Re < 10^5$) แรงต้านจะแปรผันตามกำลังสองของความเร็ว:
\begin{equation}
F_D = \frac{1}{2} C_D \rho A v^2
\end{equation}
โดยที่ $C_D$ คือสัมประสิทธิ์แรงฉุด (Drag Coefficient), $\rho$ คือความหนาแน่นของอากาศ และ $A$ คือพื้นที่หน้าตัดตั้งฉากกับทิศทางการเคลื่อนที่ ผลของแรงต้านนี้ทำให้ความเร่งของวัตถุลดลงจนเข้าสู่ศูนย์เมื่อแรงต้านสมดุลกับแรงโน้มถ่วง เกิดเป็นความเร็วปลายคงที่ (Terminal Velocity: $v_t$)

\section{ตัวอย่างการคำนวณตามกรอบวิเคราะห์ P-S-C-T}

\begin{psctexample}{2.1}{การคำนวณวิถีตกอิสระของหยดน้ำฝนและอนุภาคฝุ่นละอองในบรรยากาศ}
\textbf{1. ภาพและสถานการณ์ทางกายภาพ (Picture):}\\\\ หยดน้ำฝนรัศมี $R = 1.5\text{ mm}$ ตกจากเมฆสูง $1\,500\text{ m}$ สู่พื้นดิน ภายใต้ความหนาแน่นอากาศ $\rho = 1.225\text{ kg/m}^3$ และ $C_D = 0.45$

\vspace{4pt}
\textbf{2. กลยุทธ์และแบบจำลองทางฟิสิกส์ (Strategy):}\\\\ คำนวณมวล $m$, พื้นที่หน้าตัด $A$ แล้วแทนค่าในสมการความเร็วปลาย $v_t = \sqrt{2mg/(\rho A C_D)}$

\vspace{4pt}
\textbf{3. ขั้นตอนการคำนวณเชิงตัวเลข (Calculation):}\\\\ ปริมาตรหยดน้ำ $V = \frac{4}{3}\pi R^3 = 1.414 \times 10^{-8}\text{ m}^3$ มวล $m = \rho_w V = 1.414 \times 10^{-5}\text{ kg}$
พื้นที่หน้าตัด $A = \pi R^2 = 7.069 \times 10^{-6}\text{ m}^2$
แทนค่าหาความเร็วปลาย:
\begin{equation}
v_t = \sqrt{\frac{2(1.414 \times 10^{-5})(9.81)}{(1.225)(7.069 \times 10^{-6})(0.45)}} \approx 8.44\text{ m/s} \quad (30.4\text{ km/h})
\end{equation}
หากไม่มีแรงต้านอากาศ อัตราเร็วเมื่อถึงพื้นจะสูงถึง $v = \sqrt{2gh} = \sqrt{2(9.81)(1500)} \approx 171.5\text{ m/s}$ ($617\text{ km/h}$ ซึ่งเป็นอันตรายต่อสิ่งมีชีวิต)

\vspace{4pt}
\textbf{4. การทดสอบความสมเหตุสมผล (Test):}\\\\ ค่า $v_t \approx 8.4\text{ m/s}$ สอดคล้องอย่างแม่นยำกับผลการวัดความเร็วหยดน้ำฝนจริงด้วยดอปเปลอร์เรดาร์ตรวจอากาศ
\end{psctexample}

\section{การเชื่อมโยงสู่การวิจัยฟิสิกส์ร่วมสมัย}

\begin{researchbox}{ข้อมูลเชิงลึกจากงานวิจัย}
การจำลองการกระจายตัวของอนุภาคฝุ่นละออง PM2.5 และละอองลอยคาร์บอนแบล็คในงานวิจัยสิ่งแวดล้อม จำเป็นต้องผสานสมการแรงต้านของสโตกส์ร่วมกับ Cunningham Correction Factor เพื่อวิเคราะห์อัตราการตกตะกอนและการแขวนลอยในชั้นบรรยากาศระดับล่าง (Planetary Boundary Layer)
\end{researchbox}

\section{การฝึกแก้โจทย์ปัญหาเชิงระบบตามกรอบ C-C-A-F}

\begin{ccafsolution}{การวิเคราะห์การเคลื่อนที่แบบโพรเจกไทล์ที่มีแรงต้านตามแนวแกนคู่}
\textbf{ขั้นที่ 1: การสร้างมโนทัศน์ (Conceptualize):}\\\\ วัตถุมวล $m$ ถูกยิงด้วยความเร็วต้น $v_0$ ทำมุม $\theta$ ในระนาบ 2 มิติ ภายใต้แรงต้านกำลังสอง $\mathbf{F}_D = -c v \mathbf{v}$

\vspace{4pt}
\textbf{ขั้นที่ 2: การจำแนกประเภทโจทย์ (Categorize):}\\\\ ระบบสมการเชิงอนุพันธ์สามัญแบบไม่เชิงเส้นคู่ควบ (Non-linear Coupled ODEs)

\vspace{4pt}
\textbf{ขั้นที่ 3: การวิเคราะห์เชิงคณิตศาสตร์ (Analyze):}\\\\ สมการการเคลื่อนที่ในพิกัดคาร์ทีเซียน:
\begin{align}
m\frac{dv_x}{dt} &= -c\sqrt{v_x^2 + v_y^2}\,v_x \\
m\frac{dv_y}{dt} &= -mg - c\sqrt{v_x^2 + v_y^2}\,v_y
\end{align}
เนื่องจากสมการทั้งสองไม่สามารถหาผลเฉลยแม่นตรงในรูปฟังก์ชันมูลฐานได้ จำเป็นต้องใช้ระเบียบวิธีเชิงตัวเลข Runge-Kutta อันดับ 4 (RK4) ในการคำนวณวิถีโค้งทีละช่วงเวลา $\Delta t$

\vspace{4pt}
\textbf{ขั้นที่ 4: การสรุปและประเมินผล (Finalize):}\\\\ แบบจำลองเชิงตัวเลข RK4 เผยให้เห็นว่าระยะตกไกลสูงสุดเกิดขึ้นที่มุมต่ำกว่า $45^\circ$ เสมอเมื่อมีแรงต้านอากาศ
\end{ccafsolution}

\section{แบบฝึกหัดทบทวนและโจทย์ท้าทายประจำบท}

\begin{enumerate}[leftmargin=1.5cm, label={\textbf{\thechapter.\arabic*}}, itemsep=8pt]
    \item \textbf{โจทย์เชิงแนวคิด (Conceptual):} จงอธิบายความสำคัญทางกายภาพของกฎและสมการหลักในบทเรียนนี้ พร้อมยกตัวอย่างสถานการณ์จริงในธรรมชาติ
    \item \textbf{โจทย์การประมาณค่าเชิงฟิสิกส์ (Estimation):} จงใช้การวิเคราะห์เชิงมิติและอันดับของขนาด (Order of Magnitude) ในการประมาณค่าปริมาณสำคัญที่เกี่ยวข้อง
    \item \textbf{โจทย์วิจัยขั้นสูง (Research Challenge):} จงเขียนสคริปต์ Python จำลองพฤติกรรมของระบบตามสมการที่ได้ศึกษา พร้อมพลอตกราฟเปรียบเทียบผลลัพธ์เชิงทฤษฎี
\end{enumerate}

